\documentclass[11pt,a4paper]{article}
\usepackage[utf8]{inputenc}
\usepackage[T1]{fontenc}
\usepackage[brazil]{babel}
\usepackage{amsmath,amssymb,amsthm}
\usepackage{booktabs}
\usepackage{tikz}
\usepackage[margin=2.5cm]{geometry}
\newtheorem{proposicao}{Proposição}

\begin{document}

\section{Modelagem do problema}

\subsection{Caracterização e hipóteses}

O problema é um \emph{flowshop} de três máquinas $M_1$, $M_2$ e $M_3$, em que
cada uma das $n$ tarefas precisa ser processada nas três máquinas, nessa ordem.
Os tempos de processamento variam de tarefa para tarefa. Entre $M_1$ e $M_2$ não
pode haver espera: uma tarefa que termina em $M_1$ precisa entrar imediatamente
em $M_2$. Entre $M_2$ e $M_3$ a espera é permitida. O objetivo é minimizar o
tempo total de execução de todas as tarefas, isto é, o \emph{makespan}
$C_{\max}$. Na notação de três campos de Graham, o problema é
\[
  F3 \mid \mathit{nwt}(M_1 \to M_2),\ \mathit{prmu} \mid C_{\max}.
\]

Adotamos as hipóteses usuais de \emph{flowshop}, que o enunciado não contradiz:
\begin{enumerate}
  \item todas as tarefas estão disponíveis no instante $t=0$;
  \item os tempos de processamento são conhecidos e determinísticos, e já incluem
        eventuais tempos de preparação;
  \item não há preempção: uma operação iniciada vai até o fim sem interrupção;
  \item cada máquina processa no máximo uma tarefa por vez, e cada tarefa está em
        no máximo uma máquina por vez;
  \item o espaço de espera (\emph{buffer}) entre $M_2$ e $M_3$ é ilimitado;
  \item a restrição de não espera vale só para a passagem $M_1 \to M_2$. Uma
        tarefa pode ter o \emph{início} em $M_1$ adiado (a máquina $M_1$ fica
        ociosa) para que termine exatamente quando $M_2$ estiver livre.
\end{enumerate}

\paragraph{Por que basta considerar permutações.}
Usamos um modelo de \emph{flowshop} de permutação, em que a mesma sequência de
tarefas é usada nas três máquinas. Essa escolha não perde otimalidade:

\begin{proposicao}
Existe uma programação ótima de $F3 \mid \mathit{nwt}(M_1\to M_2) \mid C_{\max}$
em que a sequência das tarefas é a mesma nas três máquinas.
\end{proposicao}
\begin{proof}[Justificativa]
\emph{($M_1$ e $M_2$.)} Como cada tarefa entra em $M_2$ no exato instante em que
sai de $M_1$, a ordem de entrada em $M_2$ é a ordem de saída de $M_1$, que por
sua vez é a ordem de processamento em $M_1$. Logo, as duas sequências são
necessariamente iguais.

\emph{($M_2$ e $M_3$.)} Vale o argumento clássico de troca para as duas últimas
máquinas de um \emph{flowshop} com objetivo $C_{\max}$ (Conway, Maxwell e Miller,
1967): se duas tarefas consecutivas em $M_3$ aparecem na ordem inversa da de
$M_2$, trocá-las em $M_3$ não aumenta o \emph{makespan}. A troca altera apenas a
programação de $M_3$ e, portanto, não interfere na restrição de não espera entre
$M_1$ e $M_2$. Repetindo as trocas, obtém-se em $M_3$ a mesma sequência de $M_2$
sem piorar $C_{\max}$.
\end{proof}

Assim, decidir a programação equivale a escolher \emph{uma} permutação das
tarefas e calcular os instantes de término em cada máquina.

\subsection{Conjuntos e parâmetros}

\begin{itemize}
  \item $J = \{1,\dots,n\}$: conjunto de tarefas, indexado por $j$;
  \item $K = \{1,\dots,n\}$: conjunto de posições na sequência, indexado por $k$;
  \item $M = \{1,2,3\}$: conjunto de máquinas, indexado por $m$;
  \item $p_{jm} \ge 0$: tempo de processamento da tarefa $j$ na máquina $m$.
\end{itemize}

\subsection{Variáveis de decisão}

\begin{itemize}
  \item $x_{jk} \in \{0,1\}$: vale 1 se a tarefa $j$ ocupa a posição $k$ da
        sequência, e 0 caso contrário;
  \item $C_{km} \ge 0$: instante de término, na máquina $m$, da tarefa que ocupa
        a posição $k$.
\end{itemize}

A formulação é \emph{posicional}: em vez de perguntar ``quem vem antes de quem''
para cada par de tarefas, ela atribui cada tarefa a uma posição. Os tempos são
indexados pela posição, e não pela tarefa, porque as relações de precedência
nas máquinas acontecem entre posições consecutivas ($k-1$ e $k$). Para
simplificar a notação, definimos a expressão linear
\[
  P_{km} \;=\; \sum_{j \in J} p_{jm}\, x_{jk},
\]
que dá o tempo de processamento, na máquina $m$, da tarefa colocada na posição
$k$. Como apenas um $x_{jk}$ vale 1 para cada $k$, $P_{km}$ seleciona exatamente
o tempo dessa tarefa. Como os $p_{jm}$ são parâmetros, $P_{km}$ é linear nas
variáveis.

\subsection{Função objetivo}

\begin{equation}
  \min \; C_{n3} \label{eq:obj}
\end{equation}

O \emph{makespan} é o instante em que a última tarefa termina na última
máquina. Pelas restrições de capacidade de $M_3$ (abaixo), os términos em $M_3$
são crescentes com a posição, de modo que a última tarefa a terminar é a da
posição $n$. Não é preciso uma variável auxiliar $C_{\max}$.

\subsection{Restrições}

\paragraph{Atribuição.} Cada tarefa ocupa exatamente uma posição, e cada posição
recebe exatamente uma tarefa:
\begin{align}
  \sum_{k \in K} x_{jk} &= 1 && \forall\, j \in J \label{eq:atr1}\\
  \sum_{j \in J} x_{jk} &= 1 && \forall\, k \in K \label{eq:atr2}
\end{align}
Juntas, \eqref{eq:atr1}--\eqref{eq:atr2} garantem que $x$ representa uma
permutação das tarefas.

\paragraph{Máquina 1.} A primeira tarefa só pode terminar depois de processada.
As demais começam em $M_1$ depois que a anterior sai de $M_1$:
\begin{align}
  C_{11} &\ge P_{11} \label{eq:m1a}\\
  C_{k1} &\ge C_{k-1,1} + P_{k1} && \forall\, k \in K,\ k \ge 2 \label{eq:m1b}
\end{align}
A desigualdade (e não a igualdade) permite que $M_1$ fique ociosa, adiando o
início de uma tarefa. Isso é essencial: às vezes é preciso esperar para começar
em $M_1$ de modo a não terminar antes de $M_2$ estar livre.

\paragraph{Não espera entre $M_1$ e $M_2$.} A tarefa começa em $M_2$ exatamente
no instante em que termina em $M_1$. O início em $M_2$ é $C_{k2} - P_{k2}$, logo:
\begin{align}
  C_{k2} &= C_{k1} + P_{k2} && \forall\, k \in K \label{eq:nw}
\end{align}
Esta é a restrição que distingue o problema de um \emph{flowshop} comum. Com
``$\ge$'' no lugar de ``$=$'', a tarefa poderia esperar entre $M_1$ e $M_2$.

\paragraph{Capacidade da máquina 2.} A tarefa da posição $k$ só entra em $M_2$
depois que a tarefa da posição $k-1$ sai de $M_2$:
\begin{align}
  C_{k2} &\ge C_{k-1,2} + P_{k2} && \forall\, k \in K,\ k \ge 2 \label{eq:m2}
\end{align}
Substituindo \eqref{eq:nw}, esta restrição equivale a $C_{k1} \ge C_{k-1,2}$: a
tarefa só pode \emph{terminar} em $M_1$ quando $M_2$ estiver livre. É ela que,
combinada com \eqref{eq:m1b}, força o eventual atraso do início em $M_1$.

\paragraph{Máquina 3: precedência tecnológica e capacidade.} A tarefa só começa
em $M_3$ depois de sair de $M_2$ (aqui a espera é permitida, por isso
``$\ge$''), e depois que a tarefa anterior sai de $M_3$:
\begin{align}
  C_{k3} &\ge C_{k2} + P_{k3} && \forall\, k \in K \label{eq:m3a}\\
  C_{k3} &\ge C_{k-1,3} + P_{k3} && \forall\, k \in K,\ k \ge 2 \label{eq:m3b}
\end{align}

\paragraph{Domínio das variáveis.}
\begin{align}
  x_{jk} &\in \{0,1\} && \forall\, j \in J,\ k \in K \label{eq:dom1}\\
  C_{km} &\ge 0 && \forall\, k \in K,\ m \in M \label{eq:dom2}
\end{align}

\subsection{Modelo completo}

\begin{align*}
  \min\quad & C_{n3}\\
  \text{s.a.}\quad
  & \textstyle\sum_{k\in K} x_{jk} = 1 && \forall j \in J\\
  & \textstyle\sum_{j\in J} x_{jk} = 1 && \forall k \in K\\
  & C_{11} \ge \textstyle\sum_{j} p_{j1}x_{j1}\\
  & C_{k1} \ge C_{k-1,1} + \textstyle\sum_{j} p_{j1}x_{jk} && \forall k \ge 2\\
  & C_{k2} = C_{k1} + \textstyle\sum_{j} p_{j2}x_{jk} && \forall k \in K\\
  & C_{k2} \ge C_{k-1,2} + \textstyle\sum_{j} p_{j2}x_{jk} && \forall k \ge 2\\
  & C_{k3} \ge C_{k2} + \textstyle\sum_{j} p_{j3}x_{jk} && \forall k \in K\\
  & C_{k3} \ge C_{k-1,3} + \textstyle\sum_{j} p_{j3}x_{jk} && \forall k \ge 2\\
  & x_{jk} \in \{0,1\},\quad C_{km} \ge 0
\end{align*}

\paragraph{Tamanho e características.} O modelo tem $n^2$ variáveis binárias,
$3n$ variáveis contínuas e $2n + 1 + (n-1) + n + (n-1) + n + (n-1) = 7n - 2$
restrições, além do domínio. É um modelo de programação linear inteira mista
(PLIM) sem constantes ``big-$M$'', o que tende a gerar relaxações lineares mais
fortes do que as formulações disjuntivas com variáveis de precedência
$y_{ij}$. O problema é NP-difícil, já que $F3 \,||\, C_{\max}$ é NP-difícil
(Garey, Johnson e Sethi, 1976). Por isso, espera-se que o tempo de solução
cresça rapidamente com $n$, o que motiva o limite de 3600 segundos nos
experimentos.

\subsection{Exemplo ilustrativo e validação}

Considere $n=4$ tarefas com os tempos da Tabela~\ref{tab:ex}.

\begin{table}[h]
\centering
\begin{tabular}{lcccc}
\toprule
 & $T_1$ & $T_2$ & $T_3$ & $T_4$\\
\midrule
$p_{j1}$ & 3 & 5 & 2 & 4\\
$p_{j2}$ & 6 & 2 & 4 & 3\\
$p_{j3}$ & 2 & 4 & 5 & 3\\
\bottomrule
\end{tabular}
\caption{Tempos de processamento do exemplo.}
\label{tab:ex}
\end{table}

Resolvendo o modelo, obtém-se a sequência ótima $T_3, T_4, T_1, T_2$ com
$C_{\max}^* = 21$ (a sequência $T_1,T_2,T_3,T_4$ daria $C_{\max}=23$). A
Figura~\ref{fig:gantt} mostra os dois efeitos do modelo:
\begin{itemize}
  \item \textbf{Atraso proposital em $M_1$:} $T_2$ poderia começar em $M_1$ no
        instante 9, mas terminaria em 14 com $M_2$ ainda ocupada por $T_1$ (até
        15). Como não pode esperar, ela começa em 10 e termina em $M_1$
        exatamente em 15, entrando direto em $M_2$.
  \item \textbf{Espera permitida antes de $M_3$:} $T_4$ sai de $M_2$ em 9 e
        espera até 11 para entrar em $M_3$.
\end{itemize}

\begin{figure}[h]
\centering
\begin{tikzpicture}[x=0.6cm,y=0.8cm]
  \definecolor{c1}{RGB}{66,133,244}\definecolor{c2}{RGB}{219,68,55}
  \definecolor{c3}{RGB}{15,157,88}\definecolor{c4}{RGB}{244,160,0}
  \newcommand{\op}[4]{\filldraw[fill=#4!45,draw=black] (#2,#1-0.35) rectangle (#3,#1+0.35);
                      \node at ({(#2+#3)/2},#1) {\small};}
  \newcommand{\lab}[4]{\node at ({(#2+#3)/2},#1) {\footnotesize #4};}
  % M1 (y=3)
  \op{3}{0}{2}{c3}\lab{3}{0}{2}{$T_3$}
  \op{3}{2}{6}{c4}\lab{3}{2}{6}{$T_4$}
  \op{3}{6}{9}{c1}\lab{3}{6}{9}{$T_1$}
  \op{3}{10}{15}{c2}\lab{3}{10}{15}{$T_2$}
  % M2 (y=2)
  \op{2}{2}{6}{c3}\lab{2}{2}{6}{$T_3$}
  \op{2}{6}{9}{c4}\lab{2}{6}{9}{$T_4$}
  \op{2}{9}{15}{c1}\lab{2}{9}{15}{$T_1$}
  \op{2}{15}{17}{c2}\lab{2}{15}{17}{$T_2$}
  % M3 (y=1)
  \op{1}{6}{11}{c3}\lab{1}{6}{11}{$T_3$}
  \op{1}{11}{14}{c4}\lab{1}{11}{14}{$T_4$}
  \op{1}{15}{17}{c1}\lab{1}{15}{17}{$T_1$}
  \op{1}{17}{21}{c2}\lab{1}{17}{21}{$T_2$}
  % eixos
  \foreach \m/\y in {M_1/3,M_2/2,M_3/1} \node[left] at (-0.2,\y) {$\m$};
  \draw[->] (0,0.4) -- (22.5,0.4) node[right] {$t$};
  \foreach \t in {0,2,...,22} \draw (\t,0.35) -- (\t,0.45) node[below=2pt] {\scriptsize \t};
  \draw[dashed,red] (21,0.4) -- (21,3.5) node[above] {\scriptsize $C_{\max}=21$};
  \draw[dotted,thick] (9,3) -- (10,3); \node[above] at (9.5,3.35) {\scriptsize atraso};
\end{tikzpicture}
\caption{Gráfico de Gantt da solução ótima do exemplo.}
\label{fig:gantt}
\end{figure}

Para validar a formulação, o modelo foi implementado em Python (PuLP com o
\emph{solver} HiGHS) e comparado com a enumeração completa das $n!$ sequências,
usando um simulador que calcula, para cada sequência, o cronograma mais cedo
possível que respeita as restrições. Em 30 instâncias aleatórias com
$n = 3,\dots,8$ e tempos inteiros em $[1,20]$, o valor ótimo do modelo
coincidiu com o da enumeração em todos os casos.

\end{document}
